\chapter{Literature Review}
\label{chap:lr}
\chaptermark{Second Chapter Heading}


\section{Ethereum Virtual Machine}
Ethereum is a decentralized system, and at its core, the Ethereum Virtual Machine (EVM) is an essential component of the Ethereum blockchain. The EVM serves as the computational heart of the Ethereum network, enabling the execution of smart contracts and providing the foundation for a wide array of decentralized applications.

In the context of smart contract execution, it's important to mention the concept of gas limit. Gas limit is a critical aspect of Ethereum's execution model. Each operation and computation on the Ethereum Virtual Machine consumes a certain amount of gas, which is a measure of computational work. The gas limit is a cap set by the user or the entity initiating a contract execution, and it represents the maximum amount of gas they are willing to spend for that operation.

Gas limit, the EVM's stack-based architecture, and the wide set of available operations allow it to efficiently and securely for miners execute these smart contracts, granting it the status of a Turing-complete machine \cite{EthereumWhitepaper}. This level of computational flexibility empowers developers to create sophisticated programs.

\section{Existing EVM languages}
Official Ethereum documentation \cite{OfficialEthereumLanguages} enumerates only four languages that compiles to EVM: Solidity, Vyper, Yul and Fe. However, other attempts to create language for EVM are known \cite{CommunityEthereumLanguages}.

Among the available languages for EVM, Solidity stands out as the most widely adopted and utilized one \cite{SolidityWidelyUsed, SolidityDevsChallenges}. This dominance is further emphasized by data presented in \cite{SolidityVyperGithubUsage}, which reveals that there are approximately 800 times more Solidity files on Github compared to Vyper files. Consequently, it is evident that the majority of existing smart contracts on the Ethereum blockchain are written in Solidity. On one hand, the concentration of development and resources around a single language like Solidity can be viewed positively. It fosters a strong sense of community cohesion and ensures that the majority of developers are focused on refining and improving a single language, potentially enhancing its robustness and feature set. However, this dominance also brings significant challenges. The high prevalence of Solidity in the EVM ecosystem means that if a vulnerability or security flaw arises in the Solidity language, it affects a substantial portion of the smart contracts on the network. This centralized risk can be a double-edged sword, as a single point of failure in Solidity could have far-reaching consequences for the entire Ethereum blockchain. Furthermore, the Ethereum community is not without its criticisms of Solidity. Community members have voiced their concerns and challenges related to the language \cite{SolidityDevsChallenges}. These concerns range from the language's complexity to its lack of certain features, which may hinder the development of robust and secure smart contracts. So, decentralization and availability of choice is important, especially for such decentralized ecosystem as blockchain.

\subsection{Solidity}

\subsection{Vyper}

\section{Market Research}

\section{Elixir}

\section{Compiler construction}

Compilation process can be divided in 3 or 4 stages. A general overview of these stages:
Tokenization
bla bla


Parsing
bla bla


Intermediate representation (optional)
Yul Yul+ used for optimizing Simplicity used for formal proofs bla bla


Code generation

Since EVM is a relatively new virtual machine, some cases are known when other existed earlier languages gets compiled to EVM (Vyper), this approach has obvious pros: already existing lexer and parser can be utilized, there are people who are already familiar with the language and are adble to use it in development, so new people can be attracted to ethereum (maybe unrelated to this section) 

bla bla

\section{Conclusion}

% \newpage


% \begin{longtable}{c|c}
% \caption[This is the title I want to appear in the List of Tables]{Simulation Parameters} \label{table:secsimulation_params} \\
% \hline
% A & B  \\
% \hline
% \endfirsthead
% \multicolumn{2}{@{}l}{} \\
% \hline
% A & B \\
% \hline
% \endhead
% \hline
%  \textbf{Parameter} & \textbf{Value}\\
%  \hline
%  Number of vehicles & $|\mathcal{V}|$\\
%  \hline
%  Number of RSUs & $|\mathcal{U}|$\\
%  \hline
%  RSU coverage radius & 150 m\\
%  \hline
%  V2V communication radius & 30 m\\
%  \hline
%  Smart vehicle antenna height & 1.5 m\\
%  \hline
%  RSU antenna height & 25 m\\
%  \hline
%  Smart vehicle maximum speed & $v_{max}$ m/s\\
%  \hline
%  Smart vehicle minimum speed & $v_{min}$ m/s\\
%  \hline
%  Common smart vehicle cache capacities & $[50, 100, 150, 200, 250]$ mb\\
%  \hline
%  Common RSU cache capacities & $[5000,1000,1500,2000,2500]$ mb\\
%  \hline
%  Common backhaul rates & $[75, 100, 150]$ mb/s\\
%  \hline
% \end{longtable}

% \begin{figure}[hbt]
% \centering
% \includegraphics[]{figs/inno.png}
% \caption{One kernel at $x_s$ (\emph{dotted kernel}) or two kernels at
% $x_i$ and $x_j$ (\textit{left and right}) lead to the same summed estimate
% at $x_s$. This shows a figure consisting of different types of
% lines. Elements of the figure described in the caption should be set in
% italics, in parentheses, as shown in this sample caption.}
% \label{fig:secex}
% \end{figure}

% This description implies several essential properties of the task at hand:
% \begin{enumerate}
%     \item Watermark must contain all necessary information, but still, be placeable and recognizable even on smaller images. The produced watermark must be compact but have the possibility to store enough information. 
%     \item To prevent easy tampering, the watermark must be invisible to the naked eye (and, preferably, to basic image parsing tools). If malefactor does not know about the existence of watermark, they might not even try to remove it and disable it. 
% \end{enumerate}
